\section*{第 \ref{ch:b3:stem-cells} 章——干细胞、再生与衰老}

\begin{solution}{exo:b3:stem-cells:1}
\omterm{def:b3:stem-cells:stem}{自我更新}：产生至少一个仍是\omterm{def:b3:stem-cells:stem}{干细胞}的子代的分裂，
即\omterm{prop:b3:stem-cells:crypt}{隐窝}底部的 \emph{\omterm{prop:b3:stem-cells:crypt}{Lgr5}} 细胞。潜能：
细胞仍可采取的命运范围。\omterm{prop:b3:stem-cells:crypt}{隐窝}\omterm{def:b3:stem-cells:stem}{干细胞}是专能的，
可产生吸收细胞、杯状细胞、肠内分泌细胞、\omterm{prop:b3:stem-cells:crypt}{潘氏细胞}
和簇状细胞。\omterm{def:b3:stem-cells:stem}{生态位}：\omterm{prop:b3:stem-cells:crypt}{潘氏细胞}及其 Wnt、Notch 和 EGF
信号，它们使干性得以维持。\omterm{def:b3:stem-cells:stem}{过渡放大细胞}：一个已定向的子细胞，
在\omterm{prop:b3:stem-cells:crypt}{隐窝}中再分裂四到五次，然后
沿绒毛向上分化。
\end{solution}

\begin{solution}{exo:b3:stem-cells:2}
单个骨髓细胞可形成含有红细胞、粒细胞
和血小板的脾集落：一个细胞，多条谱系，
即\omterm{def:b3:stem-cells:pluripotency}{多能性}；线性剂量反应表明每个集落都来自单个细胞。
\omterm{def:b3:stem-cells:stem}{自我更新}则需要第二只小鼠：取自一个集落的细胞
注入另一只经照射的受体小鼠后又形成新的集落，
说明最初的细胞已产生出具有同等\omterm{def:b3:stem-cells:stem}{自我更新}能力的子代。
\end{solution}

\begin{solution}{exo:b3:stem-cells:3}
\emph{Oct4}、\emph{Sox2}、\emph{Klf4}、\emph{c-Myc}。成纤维细胞的
\omterm{def:b3:chromatin-epigenetics:nucleosome}{染色质}已将\omterm{def:b3:stem-cells:pluripotency}{多能性}基因包裹在\omterm{def:b3:chromatin-epigenetics:nucleosome}{异染色质}和 DNA 甲基化之下，
因此这些因子起初几乎找不到可接近的位点，
必须经过许多个\omterm{def:b3:cell-cycle-apoptosis:cdk}{细胞周期}才能招募重塑酶；
细胞还必须关闭自身程序、逃离\omterm{def:b3:cell-cycle-apoptosis:other}{衰老}，并经过一个随机的
中间状态，其中大多数细胞停滞或死亡。
只有每一步都成功的少数细胞，
才会在数周后以\omterm{def:b3:stem-cells:pluripotency}{多能性}的状态出现。
\end{solution}

\begin{solution}{exo:b3:stem-cells:4}
涡虫保留了多能的成体\omterm{def:b3:stem-cells:stem}{干细胞}，即\omterm{def:b3:stem-cells:regeneration}{新生细胞}，它们
迁移到伤口处，构建缺失的任何部分。蝾螈没有这样的细胞库：
断肢处已分化的细胞\omterm{def:b3:stem-cells:regeneration}{去分化}为\omterm{def:b3:stem-cells:stem}{祖细胞}，这些\omterm{def:b3:stem-cells:stem}{祖细胞}与原有的
组织\omterm{def:b3:stem-cells:stem}{干细胞}一起形成\omterm{def:b3:stem-cells:regeneration}{芽基}，并重新执行肢体的
发育程序。Kragl 对单一标记组织所做的移植实验表明，\omterm{def:b3:stem-cells:regeneration}{芽基}细胞
保持其来源的特性——
肌肉形成肌肉，软骨形成软骨——
因此\omterm{def:b3:stem-cells:regeneration}{芽基}是由谱系限制性\omterm{def:b3:stem-cells:stem}{祖细胞}构成的混合物，而非多能的细胞团。
\end{solution}

\begin{solution}{exo:b3:stem-cells:5}
$(11 - 5)/0.075 = 80$ 次倍增。从 \qty{8}{kb} 出发：$(8 - 5)/0.075 = 40$。
若补偿率为 \qty{60}{\%}，则净损失为 \qty{30}{bp}：$6000/30
= 200$.
\end{solution}

\begin{solution}{exo:b3:stem-cells:6}
每天 14 次\omterm{def:b3:stem-cells:stem}{干细胞}分裂产生 28 个子代；其中 14 个必须仍为\omterm{def:b3:stem-cells:stem}{干细胞}，
因此 14 个定向——每个\omterm{def:b3:stem-cells:stem}{干细胞}每天一个定向子细胞。
每个定向细胞产生 $2^{4} = 16$ 个细胞：$14\times 16 = 224$ 个绒\omterm{def:b3:sensory-systems:cochlea}{毛细胞}
每天，与观察到的 \num{300} 个数量级相当（若再多一次过渡分裂，则为 \num{448}）。
\omterm{def:b3:stem-cells:stem}{干细胞}层的贡献相当于十四个细胞的分裂，过渡放大层
贡献其余的 \num{200} 个。
\end{solution}

\begin{solution}{exo:b3:stem-cells:7}
$2\times 10^{11}/2^{20} = 1.9\times 10^{5}$ 个定向子代每天；
若有 \num{10000} 个\omterm{def:b3:stem-cells:stem}{干细胞}，则每个\omterm{def:b3:stem-cells:stem}{干细胞}每天产生 19 个定向子代，
而\omterm{def:b3:stem-cells:stem}{干细胞}每年仅分裂一次。二十次分裂的计数是
每条谱系从\omterm{def:b3:stem-cells:stem}{干细胞}到红细胞的次数，但中间\omterm{def:b3:stem-cells:stem}{祖细胞}并非一次性的：
多能与定向\omterm{def:b3:stem-cells:stem}{祖细胞}群体可通过自身分裂维持数周，
因此几乎所有分裂都发生在\omterm{def:b3:stem-cells:stem}{祖细胞}层，
\omterm{def:b3:stem-cells:stem}{干细胞}很少被调用。
\end{solution}

\begin{solution}{exo:b3:stem-cells:8}
$\mu(50) = 10^{-3}\mathrm{e}^{1.74} = 5.7\times 10^{-3}$; $\mu(70) =
10^{-3}\mathrm{e}^{3.48} = 0.032$; $\mu(90) = 10^{-3}\mathrm{e}^{5.22} =
0.19$ 每年。取 $\mu_{0}/\gamma = 0.0115$：$S(70) = \exp[-0.0115
\times 31.5] = 0.70$，$S(90) = \exp[-0.0115\times 184] = 0.12$。中位数：
$30 + \ln(61)/0.087 = 77$。$\mu_{0}$ 减半时：$\ln(121)/0.087 = 55$，
中位数为 85——即获得一个倍增时间（八年）。若将 $\gamma$ 减半至
$0.0435$：$\ln(31)/0.0435 = 79$，中位数为 109：减缓指数
所得收益是减半常数的四倍。
\end{solution}

\begin{solution}{exo:b3:stem-cells:9}
余下的三分之一须增至三倍：每个细胞需 $\log_{2}3 = 1.6$ 次分裂。
这是\omterm{def:b3:stem-cells:regeneration}{代偿性生长}，因为现存的肝叶通过
细胞原位分裂而增大；被切除的肝叶连同其导管、血管和
结构并未重建——肝脏恢复了质量和功能，
但未恢复其形态。
\end{solution}

\begin{solution}{exo:b3:stem-cells:10}
每次事件后，克隆大小 $k$ 以相等的概率变为 $k + 1$ 或 $k - 1$，
因此其期望大小永不改变；随机游走结束时，大小为 $N$（取代）或 $0$（丢失），
而期望 $N\,P + 0 = k$ 给出 $P = k/N$。单个细胞：$1/14$。
单个\omterm{def:b3:stem-cells:stem}{干细胞}中的中性突变在其\omterm{prop:b3:stem-cells:crypt}{隐窝}内固定的概率为 $1/14$，
否则会在数月内丢失——
漂变清除了十四个突变中的十三个。
\end{solution}

\begin{solution}{exo:b3:stem-cells:11}
此时随机游走有偏：克隆获得一个位置的次数多于失去一个位置的次数，
因此其期望大小增长，取代概率随 $k$ 增大而趋于必然——
若在十四个位置的\omterm{def:b3:stem-cells:stem}{生态位}中具有 \qty{10}{\%} 的优势，取代概率约为
单个细胞时 $1/14$ 的三倍，而克隆占据若干位置后，此概率还会更高。
经过数十年，\omterm{def:b3:stem-cells:stem}{干细胞}
若带有此类突变（骨髓中的 \emph{DNMT3A}、\emph{TET2}），便会取代
其\omterm{def:b3:stem-cells:stem}{生态位}。到七十岁时，五分之一的人，其大部分血液
来自同一克隆。这并非\omterm{def:b3:cancer-biology:hallmarks}{癌症}：细胞仍会分化，并遵循其层级结构。
但下一个突变现在可以在由数十亿细胞构成的克隆中出现，
而白血病的风险
因此提高到十倍。
\end{solution}

\begin{solution}{exo:b3:stem-cells:12}
中位数 $= 30 + \gamma^{-1}\ln(1 + \gamma\ln 2/\mu_{0})$：对于 $\mu_{0} =
10^{-2}$，$30 + \ln 7/0.087 = 52$；$10^{-3}$：77；$10^{-4}$：$30 +
\ln 604/0.087 = 104$。在纯 Gompertz 形式中，每次十倍的削减都
增加相同的 $\ln 10/\gamma = 26$ 年。实际中收益递减，
因为历史上的收益来自去除一个与年龄无关的项
——感染、分娩、意外，以及 Makeham 常数 $A$，在 $\mu = A
+ \mu_{0}\mathrm{e}^{\gamma t}$ 中的项——一旦 $A$ 相对于关键年龄段的指数项已经很小，
再削减它几乎不起作用；内源性的 $\mu_{0}$ 变化要小得多。
若将 $\gamma$ 降低 \qty{20}{\%}，
（至 $0.070$），则 $30 + \ln 49/0.070 = 86$：增加九年，
每个年龄段都能获益，而非靠推迟几种死因。
\end{solution}

\begin{solution}{pb:b3:stem-cells:1}
\textbf{1.} 每天 $3.5\times 10^{11}$，每秒 $4\times 10^{6}$。
\textbf{2.} 每个定向子代含 $2^{20} \approx 10^{6}$ 个细胞；
每天 $3.5\times 10^{11}/10^{6} = 3.3\times 10^{5}$ 个定向子代。
\textbf{3.} 每个\omterm{def:b3:stem-cells:stem}{干细胞}每天 $33$ 个，每年 \num{12000} 个。
\textbf{4.} \omterm{def:b3:stem-cells:stem}{祖细胞}库须自我维持：专能
与定向\omterm{def:b3:stem-cells:stem}{祖细胞}分裂数周并保持自身数量，
因此二十次分裂中几乎每一次都发生在
能\omterm{def:b3:stem-cells:stem}{自我更新}的各层级中，而\omterm{def:b3:stem-cells:stem}{干细胞}很少产生定向子代，
——量级约为每年一次。
\textbf{5.} $10 - 20\times 0.06 = \qty{8.8}{kb}$：远高于
$L_{\text{临界}}$，且红细胞不再分裂；\omterm{thm:b3:stem-cells:hayflick}{端粒}储备
只对会被重复利用的\omterm{def:b3:stem-cells:stem}{祖细胞}有意义。
\textbf{6.} 净损失 $0.3\times 60 = \qty{18}{bp}$；$6000/18 = 333$
次分裂；每年一次，需 $333$ 年——远超一生，因此
\omterm{thm:b3:stem-cells:hayflick}{端粒}并不限制\omterm{def:b3:stem-cells:stem}{干细胞}本身，限制它的是其他\omterm{def:b3:dna-repair:lesions}{损伤}。
\textbf{7.} 每天 $14$ 次分裂：$7$ 次补充丢失的\omterm{def:b3:stem-cells:stem}{干细胞}，$7$ 次
产生定向子代；每天 $7\times 2^{4} = 112$ 个绒\omterm{def:b3:sensory-systems:cochlea}{毛细胞}。
\textbf{8.} 每个细胞按每天一次计，$N^{2} = 196$ 次事件：约
\qty{200}{d}，六至七个月，这段时间内彩纸屑
\omterm{prop:b3:stem-cells:crypt}{隐窝}变为单色。
\textbf{9.} $1/14 \approx \qty{7}{\%}$；$10^{7}/14 \approx 7\times 10^{5}$
个\omterm{prop:b3:stem-cells:crypt}{隐窝}。
\textbf{10.} 漂变使每十四个突变中的十三个
在数月之内被淘汰。若进行终生的\omterm{def:b3:stem-cells:stem}{不对称分裂}，则每个突变\omterm{def:b3:stem-cells:stem}{干细胞}都将
终生保留，于是突变将在每个谱系中累积，
且永远不被清除。
\textbf{11.} 干性是\omterm{def:b3:stem-cells:stem}{生态位}中的一个位置，也是一组信号，
而非细胞的固有属性：任何到达\omterm{prop:b3:stem-cells:crypt}{潘氏细胞}
信号所及范围的细胞，均可承担这一角色。
\textbf{12.} 在某一时刻\omterm{def:b3:developmental-genetics:organiser}{诱导}\omterm{def:b3:stem-cells:stem}{干细胞}带上多色标记，
并追踪各\omterm{prop:b3:stem-cells:crypt}{隐窝}。对称漂变：\omterm{prop:b3:stem-cells:crypt}{隐窝}变为单色，
数月之内即完成。\omterm{def:b3:stem-cells:stem}{不对称分裂}：每个\omterm{prop:b3:stem-cells:crypt}{隐窝}保留全部颜色，
无限期保留，每种颜色都呈细带状沿绒毛向上延伸。
\textbf{13.} $6000/60 = 100$ 次倍增。
\textbf{14.} 约剩 $75$ 次：细胞计数的是分裂次数，而非日历
时间，因为在冷冻中度过的十年并不计入。
\textbf{15.} $(8000 - 4000)/30 = 133$ 年，自二十岁起
——即 153 岁。白细胞\omterm{thm:b3:stem-cells:hayflick}{端粒}的平均长度本身并不限制寿命；最短的
\omterm{thm:b3:stem-cells:hayflick}{端粒}及其他\omterm{def:b3:cell-cycle-apoptosis:other}{衰老}标志更为重要。
\textbf{16.} \omterm{def:b3:cancer-biology:telomerase}{端粒酶}消除了一道\omterm{def:b3:innate-immunity:innate}{屏障}——\omterm{thm:b3:stem-cells:hayflick}{复制性衰老}，
因而几乎在每种\omterm{def:b3:cancer-biology:hallmarks}{癌症}中都能找到，\omterm{def:b3:cancer-biology:hallmarks}{癌症}需要无限分裂；
但被赋予\omterm{def:b3:cancer-biology:telomerase}{端粒酶}的成纤维细胞仍保有其\omterm{def:b3:cell-cycle-apoptosis:checkpoints}{检查点}，
接触抑制及完整的 \omterm{def:b3:cancer-biology:genes}{p53}，因此仅无限分裂本身
并不足以使它们成为癌细胞。是必要条件而非充分条件。
\textbf{17.} $\mu(40) = 2.4\times 10^{-3}$，$\mu(60) = 0.014$，$\mu(80) =
0.077$，$\mu(100) = 0.44$ 每年。$S = \exp[-0.0115(\mathrm{e}^{\gamma
(t-30)} - 1)]$：40 岁时 $0.98$，60 岁时 $0.87$，80 岁时 $0.42$，
100 岁时 $0.006$。
\textbf{18.} 中位数 $30 + \ln 61/0.087 = 77$。存活比例为百分之十时，
$\mathrm{e}^{\gamma(t-30)} = 1 + \gamma\ln 10/\mu_{0} = 201$，此时 $t = 30 +
5.3/0.087 = 91$。
\textbf{19.} \qty{0.67}{mm/d}；$\log_{2}1000 = 10$ 次倍增，历时 60
天，每六天一次。
\textbf{20.} 约 $10^{8}$ 次分裂，因此每次
\omterm{def:b3:stem-cells:regeneration}{再生}产生 $0.1$ 个致癌突变——即每十个肢体中有一个。蝾螈几乎
不形成\omterm{def:b3:cancer-biology:hallmarks}{肿瘤}：它们是冷血且代谢缓慢的动物，\omterm{def:b3:cancer-biology:hallmarks}{肿瘤}抑制作用很强，
且\omterm{def:b3:stem-cells:regeneration}{芽基}细胞保持谱系限制，倾向于再分化而非继续分裂；
\omterm{def:b3:stem-cells:regeneration}{再生}组织甚至能抑制\omterm{def:b3:developmental-genetics:organiser}{诱导}产生的
\omterm{def:b3:cancer-biology:hallmarks}{肿瘤}。哺乳动物拥有更多细胞、更高的体温和代谢
率，以及更长的寿命，逃逸的克隆有更长时间生长。
\textbf{21.} 去神经支配残端，并植入释放
候选蛋白的微珠，或移植经神经处理的组织：若\omterm{def:b3:stem-cells:regeneration}{再生}得以恢复，
则表明存在一种可扩散因子。在蝾螈与
墨西哥钝口螈中鉴定出的候选因子有：前梯度蛋白 nAG、neuregulin-1 和 FGFs。
\textbf{22.} Wnt 信号在尾部高、头部低，而
断段前端的伤口通常会\omterm{def:b3:stem-cells:regeneration}{再生}出头部，因为那里的 Wnt
水平低；敲减 Wnt 抑制剂会使各处 Wnt 升高，
两处伤口均读作“后部”，于是长出两条尾巴。断段得知
自身的轴：依据肌肉中残留的梯度，伤口则
相对于该轴重新建立梯度。
\textbf{23.} 剩余的三分之一须增至三倍：每个肝细胞分裂
一次（至三分之二），一半再分裂——平均 $1.6$ 次分裂。
存活的肝叶增大至原有质量；缺失的肝叶及
其导管和血管不会重建。
\textbf{24.} 能替换磨损与受损组织的\omterm{def:b3:stem-cells:regeneration}{再生}
将减缓指数所度量的累积，降低 $\gamma$ 并
使曲线变平——蝾螈几乎不表现\omterm{def:b3:cell-cycle-apoptosis:other}{衰老}。其代价
将是在一个温暖、体型大的
长寿身体中，\omterm{def:b3:cancer-biology:hallmarks}{癌症}带来的更高常数项。
\textbf{25.} 海弗利克计数约为 $100$ 次倍增；借助\omterm{def:b3:cancer-biology:telomerase}{端粒酶}的
\omterm{def:b3:stem-cells:stem}{干细胞}约有 $330$ 次分裂，每年一次则超过三个世纪；
中位寿命为 $77$ 年。
\end{solution}
